% Part: first-order-logic
% Chapter: completeness
% Section: downward-ls

\documentclass[../../../include/open-logic-section]{subfiles}

\begin{document}

\olfileid{fol}{com}{dls}
\olsection{લેવેનહાઇમ--સ્કોલેમ પ્રમેય}

લેવેનહાઇમ--સ્કોલેમ પ્રમેય કહે છે કે જો કોઈ સિદ્ધાંતને અનંત નિદર્શ
હોય, તો તેને વધુમાં વધુ !!{denumerable} નિદર્શ પણ હોય છે. આ હકીકતનું
તાત્કાલિક પરિણામ એ છે કે પ્રથમ-ક્રમ તર્કશાસ્ત્ર !!a{structure}નું કદ
!!{nonenumerable} છે તે વ્યક્ત કરી શકતું નથી: બધી !!{nonenumerable}
!!{structure}sમાં સંતોષાતું કોઈ પણ !!{sentence} અથવા !!{sentence}sનો
ગણ કોઈક !!{enumerable} સંરચનામાં પણ સંતોષાય છે.

\begin{thm}
\ollabel{thm:downward-ls} જો~$\Gamma$ સુસંગત હોય, તો તેને
!!a{enumerable} નિદર્શ છે; એટલે કે, તે એવી !!a{structure}માં સંતોષ્ય
છે જેનું વ્યાપકક્ષેત્ર કાં તો સાન્ત અથવા !!{denumerable} છે.
\end{thm}

\begin{proof}
જો~$\Gamma$ સુસંગત હોય, તો પૂર્ણતા પ્રમેયની સાબિતીથી મળતી
!!{structure}~$\Struct M$નું વ્યાપકક્ષેત્ર~$\Domain{M}$ વિસ્તૃત ભાષા
~$\Lang L'$ના પદોના ગણથી મોટું નથી. તેથી~$\Struct M$ વધુમાં વધુ
!!{denumerable} છે.
\end{proof}

\sourcecorrection{OLCO-020}{મૂળની \texttt{downward-ls.tex},
પંક્તિ~28માં રચાયેલ નિદર્શનું વ્યાપકક્ષેત્ર મૂળ ભાષા~$\Lang L$નાં
પદોથી મોટું નથી એમ કહે છે. પરંતુ હેન્કિન રચના પહેલાં ભાષાને નવા
અચળોવાળી~$\Lang L'$ સુધી વિસ્તારે છે અને પદ-નિદર્શનાં પદો પણ
~$\Lang L'$નાં છે. અહીં યોગ્ય વિસ્તૃત ભાષા લખી છે; તેની પદસમષ્ટિ હજી
ગણનીય જ છે.}

\begin{thm}
\ollabel{noidentity-ls} જો~$\Gamma$ તાદાત્મ્ય વિનાના પ્રથમ-ક્રમ
તર્કશાસ્ત્રની ભાષામાંના !!{sentence}sનો સુસંગત ગણ હોય, તો તેને
!!a{denumerable} નિદર્શ છે; એટલે કે, તે એવી !!a{structure}માં સંતોષ્ય
છે જેનું વ્યાપકક્ષેત્ર અનંત અને !!{enumerable} છે.
\end{thm}

\begin{proof}
જો~$\Gamma$ સુસંગત હોય અને તેમાં તાદાત્મ્ય આવતું હોય એવું એકેય વાક્ય
ન હોય, તો પૂર્ણતા પ્રમેયની સાબિતીથી મળતી !!{structure}~$\Struct M$નું
વ્યાપકક્ષેત્ર~$\Domain{M}$ ભાષા~$\Lang L'$નાં પદોના ગણને અભિન્ન છે.
તેથી~$\Trm[L']$ની જેમ~$\Struct{M}$ પણ !!{denumerable} છે.
\end{proof}

\begin{ex}[સ્કોલેમનો વિરોધાભાસ]
ઝર્મેલો--ફ્રેન્કેલ ગણસિદ્ધાંત~$\Log{ZFC}$ એવું અત્યંત સમર્થ માળખું છે
જેમાં ગણોના કદ વિશેની હકીકતો સહિત લગભગ બધાં ગાણિતિક વિધાનો વ્યક્ત કરી
શકાય છે. તેથી, દાખલા તરીકે,~$\Log{ZFC}$ સાબિત કરી શકે છે કે વાસ્તવિક
સંખ્યાઓનો ગણ~$\Real$ !!{nonenumerable} છે; તે કૅન્ટૉરનું પ્રમેય સાબિત
કરી શકે છે કે કોઈ પણ ગણનો ઘાતગણ એ ગણ કરતાં મોટો છે; વગેરે. જો
~$\Log{ZFC}$ સુસંગત હોય, તો તેના બધા નિદર્શો અનંત છે. વધુમાં, તે બધા
એવાં !!{element}s ધરાવે છે જેમના વિશે સિદ્ધાંત કહે છે કે તેઓ
!!{nonenumerable} છે; જેમ કે પ્રાકૃતિક સંખ્યાઓનો ઘાતગણ અસ્તિત્વ ધરાવે
છે એ~$\Log{ZFC}$ના પ્રમેયને સાચું બનાવતો ઘટક. લેવેનહાઇમ--સ્કોલેમ
પ્રમેય મુજબ~$\Log{ZFC}$ને !!{enumerable} નિદર્શો પણ છે—એવા નિદર્શો જે
``!!{nonenumerable}'' ગણો ધરાવે છે, પરંતુ પોતે !!{enumerable} છે.
\end{ex}

\end{document}
